\section{What every additive measurement must retain}
\label{k:section}

The comparison of composition series gives more than an integer.
For each isomorphism class of simple modules $S$, let $m_S(M)$ be
the number of its occurrences in a composition series of $M$.
The common Jordan--H\"older theorem (Theorem~\ref{j:jordan-holder})
makes this well-defined. Only finitely many of these numbers are
nonzero for a fixed finite-length module $M$. The universal monoid
statement was proved in Corollary~\ref{j:universal-monoid}. We now
construct its group-valued recipient explicitly for modules, including
the size and closure checks needed to use it.

\subsection{Splicing gives additivity before we name a group}

The last construction in Section~\ref{m:section} spliced composition
series across a short exact sequence. Uniqueness now turns that
particular construction into an intrinsic assertion:
\begin{equation}\label{k:multiplicity-additivity}
0\to A\to E\to B\to0
\quad\Longrightarrow\quad
m_S(E)=m_S(A)+m_S(B)\quad\hbox{for every simple }S.
\end{equation}
The existence arguments there also show that the category of
finite-length $R$-modules is closed under submodules, quotients and
extensions. Every object subsequently used in such a sequence stays
inside that category.

An additive measurement is an assignment $a(M)$ in an abelian group
$D$, constant on isomorphic modules, such that
$a(E)=a(A)+a(B)$ for every short exact sequence of finite-length
modules. The zero sequence forces $a(0)=0$. Repeatedly using the
short exact sequences of a composition series gives
\begin{equation}\label{k:all-measurements}
 a(M)=\sum_S m_S(M)a(S).
\end{equation}
Thus a choice of values on simple modules determines every additive
measurement. Conversely any such choice defines an additive measurement
by this formula and \eqref{k:multiplicity-additivity}. This already
answers the universal question in concrete terms.

\subsection{Constructing the universal recipient}

For completeness we construct the group that receives every additive
measurement. Take one formal generator $e_M$ for each isomorphism
class of finite-length $R$-modules. These classes form a set: every
simple module is cyclic, since the map $R\to S$, $r\mapsto rx$, has
nonzero image for $x\ne0$; inductively every finite-length module is
finitely generated by generators of a submodule and lifts of generators
of a quotient. Such a module is a quotient of some $R^n$, and the
submodules of all these $R^n$ form a set.

Let $F$ be the abelian group of finitely supported integer combinations
of the symbols $e_M$, with coefficientwise addition. Let $J$ be the
subgroup generated by
\[
 e_E-e_A-e_B
 \quad\text{for all }0\to A\to E\to B\to0.
\]
Define $K_0(R\text{-}\mathrm{fl})=F/J$, and write $[M]$ for the
image of $e_M$. Here $R\text{-}\mathrm{fl}$ means finite-length
left $R$-modules. We specify that category because enlarging it can
introduce new exact sequences and therefore new relations.

Given an additive measurement $a$ into $D$, the formula
$\sum n_Me_M\mapsto\sum n_Ma(M)$ is a homomorphism $F\to D$.
It kills $J$ exactly by additivity, so the quotient universal property
gives a unique homomorphism $\bar a:K_0\to D$ with $\bar a([M])=a(M)$.
Conversely such a homomorphism defines an additive measurement.
This is the universal property of the Grothendieck group.

\begin{theorem}\label{k:free-simples}
The homomorphism
\[
 \bigoplus_{[S]\text{ simple}}\mathbb Z e_S\longrightarrow
 K_0(R\text{-}\mathrm{fl}),\qquad e_S\longmapsto[S]
\]
is an isomorphism. Its inverse sends $[M]$ to
$\sum_Sm_S(M)e_S$.
\end{theorem}
\begin{proof}
Additivity of every $m_S$ shows that the proposed inverse kills each
defining relation and is therefore well-defined. A simple module has
one factor, itself, so the composite on each $e_S$ is the identity.
For the other composite, the defining relations applied to a composition
series give $[M]=\sum_Sm_S(M)[S]$. Since the classes $[M]$ generate
$K_0$, this composite too is the identity.
\end{proof}

The direct sum in the theorem means finite support, even if there are
infinitely many simple isomorphism classes. The theorem proves exactly
which information all additive measurements retain. It says nothing
about measurements that fail additivity on short exact sequences.

Length is the homomorphism sending every $[S]$ to $1$. For
finite-dimensional vector spaces there is one simple class $[k]$,
so $K_0\cong\mathbb Z$ and $[V]=(\dim_kV)[k]$.
For finite-dimensional modules over a $k$-algebra, dimension instead
assigns the weight $\dim_k S$ to a simple class. To use the same proof
in that category, observe that submodules and quotients are finite-dimensional
by Lemma~\ref{m:bases}, extensions have a basis made from a submodule
basis and quotient lifts, and all these modules have composition series
by Section~\ref{m:length-foundations}.

\subsection{Why the chosen category matters}

Let $S=k[t]/(t)$. In the category of finite-length $k[t]$-modules,
$[S]$ is a nonzero generator of the free group in the theorem.
If we instead impose exact-sequence relations for all finitely generated
$k[t]$-modules, the sequence
\[
0\longrightarrow k[t]\xrightarrow{\,t\,}k[t]
\longrightarrow S\longrightarrow0
\]
forces $[S]=0$. Multiplication by $t$ is injective by inspection of
coefficients, and its image is exactly the polynomials with zero
constant term, proving exactness. The module $k[t]$ does not have
finite length: the strictly descending submodules $(t^n)$ would each
force a strictly smaller nonnegative length if it did, by
\eqref{k:multiplicity-additivity}. Thus this new relation was not
available in the previous category.

The invariant depends on which objects and exact sequences we permit.
Keeping that context visible prevents a numerical formula from
silently changing its meaning.

\subsection{A short bridge to Euler characteristic}
\label{k:euler}

Let $C^\bullet$ be a bounded complex of finite-length modules: it
has maps $d^i:C^i\to C^{i+1}$ with $d^{i+1}d^i=0$, and only
finitely many $C^i$ are nonzero. Put
\[
 Z^i=\Ker d^i,\quad B^i=\im d^{i-1},\quad H^i=Z^i/B^i.
\]
The complex condition gives $B^i\subseteq Z^i$. Kernels, images and
quotients are finite-length by the results already proved. The two
short exact sequences
\[
0\to Z^i\to C^i\xrightarrow{d^i}B^{i+1}\to0,
\qquad 0\to B^i\to Z^i\to H^i\to0
\]
give $[C^i]=[B^i]+[H^i]+[B^{i+1}]$. Alternating addition cancels
every boundary class with its next occurrence. Boundedness makes all
the sums finite, and proves
\begin{equation}\label{k:euler-formula}
 \sum_i(-1)^i[C^i]=\sum_i(-1)^i[H^i].
\end{equation}
Applying any additive measurement gives its numerical version.
In particular, vector-space dimension gives the usual Euler
characteristic formula.

This identity follows directly from the exact-sequence relations.
We have introduced it only after the earlier comparison. In
particular, we did not infer Jordan--H\"older from
an equality in a Grothendieck group whose simple generators had
already been justified by Jordan--H\"older. The preceding spectral
argument compared actual simple modules and actual differentials.

\subsection{The information that remains to be recovered}

If $0\to A\to E\to B\to0$ is exact, its class is always
$[A]+[B]$, whether or not the sequence splits. Every additive
measurement therefore agrees on $E$ and $A\oplus B$.
This is an exact limitation of additivity. The next chapter retains
the maps of the extension, and measures how the middle object can
depend on them.
